% !TeX program = xelatex
% ============================================================
% interim_report/interim_report.tex (EN) — "Interim report" of the research
% course project, SE programme — the CP1 deliverable.
%
% Structure per §2.3.2.4 of the Practice Programme (CP1 reporting for a
% research project): title page (with UDC) → contents → key terms,
% definitions and abbreviations → introduction (relevance, object and
% subject, methods, aim and objectives; novelty and validity of the
% expected results, their theoretical significance and practical value) →
% review and analysis of sources, selection of methods/algorithms/models →
% list of references (GOST) → Appendix — work schedule with dates.
%
% This is NOT the final Report (GOST 7.32-2017): no abstract, experimental
% part or conclusion here — those belong to the "Report" document for the
% defence (CP3). Submitted at CP1 without signatures (pass/fail). The
% interim report is not resubmitted at CP2 (CP2 = source-code link +
% updated project description).
%
% How to use: replace the yellow \hseFill{…} stubs with your text; set your
% personal data, supervisor, group, year and UDC in the project settings.
% ============================================================
\documentclass[12pt,a4paper]{extreport}
\input{preamble}
\begin{document}
\input{title}

% ── Contents ──
\tableofcontents
\clearpage

% ============================================================
\chapter*{KEY TERMS, DEFINITIONS AND ABBREVIATIONS}
\addcontentsline{toc}{chapter}{Key terms, definitions and abbreviations}
% ============================================================
The following terms, definitions and abbreviations are used in this report.

\begin{description}[leftmargin=0pt, labelindent=0pt, font=\bfseries, style=nextline]
  \item[\hseFill{Term 1}] \hseFill{definition}.
  \item[\hseFill{Term 2}] \hseFill{definition}.
  \item[\hseFill{ABBR}] \hseFill{expansion of the abbreviation}.
\end{description}

\clearpage
% ============================================================
\chapter*{INTRODUCTION}
\addcontentsline{toc}{chapter}{Introduction}
% ============================================================
% §2.3.2.4: relevance, object and subject, methods, aim and objectives;
% novelty and validity of the expected results, their significance.
\textbf{Relevance.} \hseFill{what problem is being solved, who it affects and
why the existing approaches are insufficient}~\cite{gost732}.

\textbf{Object of research} — \hseFill{what is studied as a whole}.

\textbf{Subject of research} — \hseFill{the specific aspect of the object}.

\textbf{Research methods:} \hseFill{which methods will be used to solve the
objectives (literature analysis, mathematical modelling, computational
experiment, empirical evaluation, etc.)}.

\textbf{Aim of the work} — \hseFill{a single statement of the scientific
result to be achieved}.

\textbf{Objectives:}
\begin{enumerate}[label=\arabic*), leftmargin=1.25cm]
  \item \hseFill{review the existing approaches};
  \item \hseFill{propose and formalise a method (model, algorithm)};
  \item \hseFill{implement a software tool and run an experiment};
  \item \hseFill{analyse the results and draw conclusions}.
\end{enumerate}

\textbf{Novelty and validity of the expected results.}
\hseFill{what is novel about the expected results and how their validity
will be ensured (theoretical justification, reproducible experiment,
comparison with baseline methods)}.

\textbf{Theoretical significance and practical value.}
\hseFill{what the expected results contribute to theory and where they can
be applied in practice}.

\clearpage
% ============================================================
\chapter{REVIEW AND ANALYSIS OF SOURCES, SELECTION OF METHODS}
% ============================================================
\section{Review and analysis of sources}
Give an analytical review of the domain and the existing approaches with
citations~\cite{article, webresource}. Every source in the reference list
must be cited in the text.

\section{Comparative analysis of existing approaches}
Compare the known methods against meaningful criteria; the comparison is
given in Table~\ref{tab:approaches}.

\begin{table}[H]
\caption{Comparison of existing approaches}
\label{tab:approaches}
\centering
\begin{tabular}{|>{\raggedright\arraybackslash}p{5cm}|c|c|c|}
\hline
\textbf{Criterion} & \textbf{Approach~A} & \textbf{Approach~B} & \textbf{Proposed} \\
\hline
Criterion~1 & \cmpplus  & \cmpminus  & \cmpplus \\
\hline
Criterion~2 & \cmpminus & \cmpplanned & \cmpplus \\
\hline
\end{tabular}
\end{table}

\section{Selection of methods, algorithms and models}
Justify the selection of methods, algorithms and models for solving the
stated objectives: \hseFill{which methods are selected and why they fit
this particular problem}.

\clearpage
% ── References. The Practice Programme requires the reference list of any
%    interim documentation to follow the GOST bibliographic standards
%    (GOST R 7.0.100-2018 / 7.0.5-2008); IEEE-style citations apply to the
%    English-language final texts, not to this GOST-formatted list. ──
\phantomsection
\addcontentsline{toc}{chapter}{References}
\renewcommand{\bibname}{References}
\begin{thebibliography}{99}
  \bibitem{article}
  Surname, N.~P. Title of the article / N.~P.~Surname // Journal Title. --- 2024. --- No.~1. --- P.~10--20.
  \bibitem{webresource}
  Resource title [Electronic resource]. --- URL: \url{https://example.org} (accessed: 01.01.2026).
  \bibitem{gost732}
  GOST 7.32--2017. The research report. Structure and rules of presentation. --- Moscow: Standartinform, 2017. --- 27~p.
\end{thebibliography}

\clearpage
% ── Appendix — work schedule with dates (§2.3.2.4) ──
\appendix
\titleformat{\chapter}[hang]
  {\normalfont\large\bfseries\centering}{Appendix~\thechapter}{1ex}{}
\chapter{Work schedule}
\begin{longtable}{|>{\raggedright\arraybackslash}p{3.2cm}|>{\raggedright\arraybackslash}p{5.2cm}|>{\raggedright\arraybackslash}p{3.6cm}|>{\raggedright\arraybackslash}p{2.6cm}|}
\hline
\textbf{Stage} & \textbf{Work description} & \textbf{Expected result} & \textbf{Dates} \\ \hline
\endhead
\hseFill{stage 1} & \hseFill{what is done} & \hseFill{result} & \hseFill{DD.MM--DD.MM} \\ \hline
\hseFill{stage 2} & \hseFill{what is done} & \hseFill{result} & \hseFill{DD.MM--DD.MM} \\ \hline
\hseFill{stage 3} & \hseFill{what is done} & \hseFill{result} & \hseFill{DD.MM--DD.MM} \\ \hline
\end{longtable}


\end{document}
